% Einheit 2 von 4 – Punktprobe und Punkte auf der Geraden (bank/geraden/e2.jsonl, zusammenbau v0.1)
%% TODO Zeitmarke Einheit 2: Marken-Zeile nicht lesbar
%% TODO Zweigzeile Teil 1: was hier gelernt wird (Satz in Schülersprache aus dem Einheitstitel) – Einheit 2
\einheitenkopf[e2]{Einheit 2 von 4 · Punktprobe und Punkte auf der Geraden}
\zweigzeile{Abitur GK}
\setcounter{aufgabe}{12}

%% TODO Titel: Ich-kann-Satz fehlt in der Bank (Platzhalter: Kettenname) – Nr. 13
%% TODO Anweisung: ein Satz über der ersten Teilaufgabe fehlt in der Bank – Nr. 13
\begin{aufgabe}{Punktprobe und Punkte auf der Geraden}
%% TODO \rechenplatz{n}: Schrittzahl des Grundfalls fehlt in der Bank (nur bei fester Schreibform, 2.3 b)
\begin{teile}
\teil Für $g\colon \vec x = (2 | 5 | 1) + t \cdot (0 | 3 | 4)$ und $P(2 | 11 | 9)$: Aus welcher Koordinate lässt sich $t$ bestimmen? \kreuz{nur aus der ersten} \\ \kreuz{aus der zweiten oder der dritten} \\ \kreuz{aus keiner}
\teil Prüfe, ob $P(7 | 5 | 8)$ auf $g\colon \vec x = (1 | 2 | -1) + r \cdot (2 | 1 | 3)$ liegt.
\teil Gegeben sind $g\colon \vec x = (5 | -1 | 2) + t \cdot (3 | 0 | -2)$, $t \in \mathbb{R}$, und $A(2 | 3 | 4)$. Begründe, dass $A$ nicht auf $g$ liegt. \hfill (Abitur 2023 GK)
\teil Der Punkt $B(-2 | y_B | z_B)$ liegt auf $g\colon \vec x = (1 | 3 | 2) + r \cdot (-1 | 2 | -3)$. Bestimme $y_B$ und $z_B$. \hfill (Abitur 2019 GK) \\ $y_B$ = \leerfeld , $z_B$ = \leerfeld
\teil Die Gerade $g$ verläuft durch $A(3 | 1 | 2)$ mit dem Richtungsvektor $\vec u = (2 | -1 | 2)$. Bestimme die Koordinaten eines Punkts auf $g$, der von $A$ den Abstand 9 hat. \hfill (Abitur 2019 GK)
\teil Eine Hecke verläuft geradlinig von $E(8 | 0 | 0)$ nach $F(0 | 8 | 0)$ (1 LE = 1 m). Weise nach, dass der Schattenpunkt $T(2 | 6 | 0)$ auf der Strecke $EF$ liegt. \hfill (Abitur 2018 GK)
\teil Weise nach, dass die Punkte $A(2 | -1 | 3)$, $B(4 | 0 | 1)$ und $C(8 | 2 | -3)$ auf einer Geraden liegen. \hfill (Abitur 2020 GK)
\teil Ein Sessellift führt geradlinig von $A(1 | 20 | 2)$ nach $E(7 | 2 | 11)$; die Talstation liegt im Ursprung, die $x_1x_2$-Ebene ist waagerecht, 1 LE entspricht 5 m. Wie hoch über der Talstation ist der Lift an der Stelle, die 84 m vom Anfang $A$ entfernt ist? \hfill (Abitur 2023 GK) \\ \leerfeld[m]
\teil Ein Prisma hat die Ecken $A(5 | 0 | 0)$ und $G(1 | 4 | 6)$. Gib eine Gleichung der Geraden durch $A$ und $G$ an und weise nach, dass der Punkt $S(3 | 2 | 5)$ nicht auf ihr liegt. \hfill (Abitur 2026 GK)
\teil Gegeben ist die Geradenschar $g_k\colon \vec x = (2 + 4k | -2k | 3 + 6k) + r \cdot (-2 | 1 | -3)$, $r \in \mathbb{R}$, für jede reelle Zahl $k$. Weise nach, dass der Koordinatenursprung für keinen Wert von $k$ auf $g_k$ liegt. \hfill (Abitur 2025 LK)
\end{teile}
\end{aufgabe}

%% TODO Titel: Ich-kann-Satz fehlt in der Bank (Platzhalter: Kettenname) – Nr. 14
%% TODO Anweisung: ein Satz über der ersten Teilaufgabe fehlt in der Bank – Nr. 14
\begin{aufgabe}{Punkt auf einer Geraden mit vorgegebenem Abstand zu einem festen Punkt über eine quadratische Gleichung bestimmen}
\begin{teile}
\teil An einem Gartenpavillon steht ein Pfosten senkrecht; sein oberes Ende ist $E(2 | 1 | 3)$, der waagerechte Balken führt von $E$ nach $F(2{,}9 | 2{,}2 | 3)$ (1 LE = 1 m). Eine $1{,}5$ m lange Strebe ist am Pfosten in $I(2 | 1 | 1{,}8)$ befestigt, ihr anderes Ende am Balken $EF$. In welchem Verhältnis teilt der obere Befestigungspunkt den Balken? \hfill (Abitur 2017 LK)
\end{teile}
\end{aufgabe}

%% TODO Titel: Ich-kann-Satz fehlt in der Bank (Platzhalter: Kettenname) – Nr. 15
%% TODO Anweisung: ein Satz über der ersten Teilaufgabe fehlt in der Bank – Nr. 15
\begin{aufgabe}{Punktprobe und Punkte auf der Geraden – Fehler finden, Begründen, Anwendung}
\begin{teile}
\teil Finde den Fehler. Prüfe, ob $P(5 | 1 | 7)$ auf $g\colon \vec x = (1 | 3 | 2) + r \cdot (2 | -1 | 1)$ liegt. Mias Rechnung: erste Koordinate $1 + 2r = 5$, also $r = 2$; zweite Koordinate $3 - 2 = 1$ ✓. Also liegt $P$ auf $g$.
\teil Begründe, warum bei der Punktprobe ein einziger Widerspruch genügt, um zu zeigen, dass der Punkt nicht auf der Geraden liegt.
\teil Ein Heißluftballon steigt geradlinig vom Startplatz im Ursprung zum Punkt $B(300 | 400 | 1\,200)$ (Angaben in m, $x_3$ ist die Höhe). Ab welchem Punkt seiner Bahn ist er höher als 900 m?
\end{teile}
\end{aufgabe}
